% LTeX: language=en-US

\documentclass[handout,usepdftitle=false,xcolor=dvipsnames,aspectratio=169]{beamer}
%\documentclass[beamer,usepdftitle=false,xcolor=dvipsnames,aspectratio=169]{beamer}
%\documentclass[trans,notes=show,usepdftitle=false,aspectratio=169]{beamer}


\usepackage[english]{babel}
\usepackage{iftex}
\ifPDFTeX
\usepackage[T1]{fontenc}
\usepackage[utf8]{inputenc}
\usepackage{lmodern}
\else
\usepackage{fontspec}
\defaultfontfeatures{Ligatures=TeX}
\setmainfont{Latin Modern Roman}
\setsansfont{Latin Modern Sans}
\setmonofont{Latin Modern Mono}
\usepackage{newunicodechar}
\newfontfamily\unicodeSymbolFont{Segoe UI Symbol}[
  BoldFont=Segoe UI Symbol,
  ItalicFont=Segoe UI Symbol,
  BoldItalicFont=Segoe UI Symbol
]
\newunicodechar{⟂}{{\unicodeSymbolFont ⟂}}
\fi
\usepackage{amsmath,amsfonts}
\usepackage{booktabs,multirow,xcolor,colortbl}
\usepackage{bm}%to use a more powerful version of \boldsymbol
\usepackage{fvextra}
\usepackage{csquotes}
\usepackage{textcomp}
\ifPDFTeX
\DeclareUnicodeCharacter{27C2}{\ensuremath{\perp}}
\fi

\newcommand{\AllSample}{ \mathcal Y^T }

\usepackage[citestyle=authoryear-comp,bibstyle=../Preamble/mystyle_lowercase,maxbibnames=5,minbibnames=1,maxnames=2, minnames=1,datezeros=false,date=long,isbn=false,url=false,backend=biber,useprefix=true]{biblatex}
\addglobalbib{../Preamble/Macro_database_biblatex.bib}	

\usepackage{upquote}
\usepackage{ragged2e}
\usepackage{microtype}

\usepackage{multimedia}
\usepackage{graphicx}
\usepackage[export]{adjustbox}
\usepackage{epstopdf}
\usepackage[labelfont=bf,format=hang, figurename=Fig.]{caption}
\usepackage[subrefformat=parens,labelformat=empty]{subcaption}
\usepackage{pgfpages} % for multiple-screen presentations

\usepackage{array}

\newcolumntype{L}[1]{>{\RaggedRight\let\newline\\\arraybackslash\hspace{0pt}}m{#1}}
\newcolumntype{C}[1]{>{\centering\let\newline\\\arraybackslash\hspace{0pt}}m{#1}}
\newcolumntype{R}[1]{>{\raggedleft\let\newline\\\arraybackslash\hspace{0pt}}m{#1}}
%https://tex.stackexchange.com/questions/12703/how-to-create-fixed-width-table-columns-with-text-raggedright-centered-raggedlef


%%% use serif font for math
\usefonttheme[onlymath]{serif}

%%% define blue highlighting
\newcommand{\highlight}[1]{{\color{blue}{#1}}}

%%%% define command for setting sources

\newcommand{\Source}{}

%%%
\beamertemplatenavigationsymbolsempty


%%% define themes
\usetheme{Boadilla}
\usecolortheme{rose}
\usecolortheme[named=blue]{structure}

\newcommand\textbox[1]{%
  \parbox{.333\textwidth}{#1}%
}

\setbeamertemplate{sections/subsections in toc}[sections numbered]


%%%% set coloring in  foot and head
\setbeamercolor{alerted text}{fg=red!80!black}
\setbeamertemplate{enumerate items}[default] % numbers enumerate items without encircling them
\setbeamertemplate{itemize items}[ball] % numbers enumerate items without encircling them
\setbeamercolor{section in head/foot}{fg=black, bg=white}
\setbeamercolor{subsection in head/foot}{fg=black, bg=white}
\setbeamercolor{title}{fg=black}
\setbeamercolor{subtitle}{fg=blue}

\setbeamertemplate{caption}[numbered]

%%% set overlays as default
\beamerdefaultoverlayspecification{<+->}

%%%% settings for output mode

\mode<beamer>
{
%\setbeameroption{show notes on second screen=right}
}

\mode<trans>
{\renewcommand{\highlight}[1]{{\textbf{#1}}}}


\mode<handout>
{
\renewcommand{\highlight}[1]{{\textbf{#1}}}
%\pgfpagesuselayout{4 on 1}[a4paper,border shrink = 5mm,landscape]
%\setbeamercolor{background canvas}{bg=black!5}
}


%%% define Exercise counter
\newcounter{Exercise}
\resetcounteronoverlays{Exercise}
\setcounter{Exercise}{0}
\newcommand{\Exercise}{\refstepcounter{Exercise}\emph{\structure{Übung \theExercise}: }}

%%%% define environments for theorems etc and make them numbered
\setbeamertemplate{theorems}[numbered]

\newtheorem{assum}{Assumption}
\setbeamertemplate{assum}[numbered]

\newtheorem{defin}{Definition}
\setbeamertemplate{defin}[numbered]

\newtheorem{propos}{Proposition}
\setbeamertemplate{propos}[numbered]


%%%% Make Parts show up in TOC
\makeatletter
\AtBeginPart{%
  \addtocontents{toc}{\protect\beamer@partintoc{\the\c@part}{\beamer@partnameshort}{\the\c@page}}%
}
%% number, shortname, page.
\providecommand\beamer@partintoc[3]{%
  \ifnum\c@tocdepth=-1\relax
    % requesting onlyparts.
    \makebox[6em]{#1.} #2
    \par
  \fi
}
\define@key{beamertoc}{onlyparts}[]{%
  \c@tocdepth=-1\relax
}
\makeatother%

%%%% fix bug in beamer package with note itemize
%%% http://tex.stackexchange.com/questions/1480/changing-the-textwidth-of-the-notes-in-beamer-repost-from-so
\makeatletter
\defbeamertemplate{note page}{infolines}
{%
  \vskip.75em
  \nointerlineskip
  \begin{minipage}{\textwidth} % this is an addition
  \footnotesize \insertnote
  \end{minipage}               % this is an addition
}
\makeatother

\setbeamertemplate{note page}[infolines]


\newcommand{\backupbegin}{
   \newcounter{framenumberappendix}
   \setcounter{framenumberappendix}{\value{framenumber}}
}
\newcommand{\backupend}{
   \addtocounter{framenumberappendix}{-\value{framenumber}}
   \addtocounter{framenumber}{\value{framenumberappendix}}
}

\hypersetup{bookmarksnumbered=true,naturalnames=true,pdfhighlight=/N,citebordercolor={1 1
1},linkbordercolor={1 1 1},colorlinks=true,anchorcolor=black,linkcolor=blue,citecolor=green,urlcolor=green,
pdfpagemode=UseOutlines,breaklinks=true,
pdfkeywords=Dynare, pdfsubject=Dynare,pdfstartpage=1} %pdfpagemode=fullpage

\usepackage[copyright]{ccicons}

\usepackage{tikz,pgfplots}
\pgfplotsset{compat=1.18}
\usetikzlibrary{arrows,calc,decorations.pathreplacing,decorations.shapes}
\pgfdeclarelayer{background}
\pgfdeclarelayer{foreground}
\pgfsetlayers{background,main,foreground}

\usepackage{relsize}
\newcommand\LM{\ensuremath{\mathit{LM}}}
\newcommand\IS{\ensuremath{\mathit{IS}}}
\DeclareMathOperator{\std}{std}


\setbeamertemplate{footline}{
\begin{beamercolorbox}{subsection in head/foot}
\insertsectionnavigationhorizontal{.95\textwidth}{\hskip0pt
\hfill}{\hfill \insertframenumber /\inserttotalframenumber}
\end{beamercolorbox}
}

\newenvironment{witemize}{\itemize\addtolength{\itemsep}{8pt}}{\enditemize}
\newenvironment{wenumerate}{\enumerate\addtolength{\itemsep}{8pt}}{\endenumerate}

\usepackage[]{minted}
\usepackage{tcolorbox}
\usepackage{etoolbox}
\BeforeBeginEnvironment{minted}{\begin{tcolorbox}[boxsep=0mm]}%
\AfterEndEnvironment{minted}{\end{tcolorbox}}%

\AtBeginEnvironment{minted}{%
  \renewcommand{\fcolorbox}[4][]{#4}}  

\usepackage{smartdiagram}

\usepackage{upquote}


\begin{document}


\title[]{\textbf{Advanced Dynare\\}
  \medskip{Chapter XIV}\\
  \emph{Optimal Policy and Welfare Assessment}
}

\author{Johannes Pfeifer}
\input{../date_place.txt}

\AtBeginSection[]
{
  \begin{frame}
    \frametitle{Outline}
    \tableofcontents[currentsection]
  \end{frame}
}


\maketitle
%\include{Intro}
\begin{frame}
  \frametitle{Outline}
  \tableofcontents%[onlyparts]
\end{frame}



\section{Introduction}

\begin{frame}{Motivation}
  \begin{witemize}
    \item Policy-makers are often interested in implementing some form of optimal policy
    \item Typical problem structure involves:
    \begin{enumerate}
      \item optimizing some welfare of loss criterion,
      \item taking into account private sector reactions to policy changes
      \item while being constrained by technologies, institutions, etc.
    \end{enumerate}
    \item Welfare analysis allows
    \begin{witemize}
      \item model-consistent comparison of policies and institutions
      \item preventing socially inefficient equilibria
    \end{witemize}
  \end{witemize}
\end{frame}

\begin{frame}{Rules vs. discretion}
  \begin{witemize}
    \item In dynamic models, private sector expectations about government behavior matter for private decision-making \parencite[e.g.][]{Lucas1976}
    \item However: government may have incentive to renege on earlier promises made to private sector (time-consistency problem of \textcite{KydlandPrescott1977})
    \item If private sector anticipates reneging, it will adjust its behavior accordingly\\
    $\rightarrow$ policy-tradeoffs like sacrifice ratio can become more expensive
    \item Solution may be credibly committing to future policy actions
  \end{witemize}
\end{frame}

\begin{frame}{Two important conceptual distinctions}
  \begin{witemize}
    \item First important conceptual distinction between
    \begin{enumerate}
      \item \highlight{discretionary policy}\\
            $\rightarrow$ policy-maker re-optimizes period-by-period
      \item optimal policy under \highlight{commitment}\\
            $\rightarrow$ announce a policy plan and stick to it no matter what
    \end{enumerate}
    \item Policy under commitment in turn can be classified in two kinds of policies
    \begin{enumerate}
      \item fully state-contingent, i.e. \highlight{Ramsey} optimal policy\\
            $\rightarrow$ first-best policy that can optimally react to all state variables including shocks
      \item \highlight{optimal simple rules (OSR)}\\
            $\rightarrow$ policy rule restricted to a specific functional form (e.g.\ Taylor rule or fiscal feedback)
    \end{enumerate}
    \item OSR is easier to communicate and implement, but may only be associated with a small welfare loss compared to Ramsey policy \parencite[e.g.][]{SGU2007}
  \end{witemize}
\end{frame}

\section{Optimal policy under commitment}

\subsection{Theory}

\begin{frame}{General case}

  \begin{witemize}
    \item  The generic Ramsey problem can be written as
          \begin{equation}
            \max_{\left\{y_t\right\}_{t=0}^\infty} \mathbb{E}_0 \sum_{t=0}^\infty\beta^{t}U(y_t)  \: ,
          \end{equation}
          where $U()$ is the planner objective and $\beta$ the planner discount factor
    \item Optimization is conducted subject to the private sector equilibrium conditions:
          \begin{equation}
            \mathbb{E}_{t}f(y_{t+1}, y_{t}, y_{t-1}, \varepsilon_{t})=0,~\forall \:  t \geq 0 \: ,
          \end{equation}
          taking $y_{-1}$ as given and imposing suitable transversality conditions, where
          \begin{itemize}
            \item $y_t\in \mathbb{R}^n$: endogenous variables
            \item $\varepsilon_t\in \mathbb{R}^p$: stochastic innovations
            \item $f~:~\mathbb{R}^{3n+p} \rightarrow \mathbb{R}^m$
          \end{itemize}
    \item There are $n-m$ free policy instruments
  \end{witemize}
\end{frame}

\begin{frame}{The Ramsey planner's Lagrangian}
  \begin{witemize}
    \item Denoting with $\mu_t \in \mathbb{R}^m$ the vector of \emph{discounted} Lagrange multipliers, the Lagrangian of the Ramsey problem is given by
    \begin{equation}
      \mathcal{L} = \mathbb{E}_0 \sum_{t=0}^\infty\beta^{t} \left( U(y_t) - \mu_t' f(y_{t+1}, y_{t}, y_{t-1}, \varepsilon_{t}) \right)
    \end{equation}
    \item Dynare will automatically construct this Lagrangian internally and compute the necessary first order conditions if a \texttt{ramsey\_model} statement is present
    \item Resulting first-order conditions can be inspected using Dynare's \LaTeX\ export functionality (\texttt{write\_latex\_dynamic\_model})
  \end{witemize}
\end{frame}

\begin{frame}{First order necessary conditions}
  \begin{itemize}
    \item The first order necessary conditions (FOC) for an optimum are given by
          \begin{align}
            0= \frac{\partial \mathcal{L}}{\partial \mu_t} = \mathbb{E}_0 & f(y_{t+1}, y_{t}, y_{t-1}, \varepsilon_{t}),~t\geq 0                                                                                                                                                                                                                                      \\
            0=\frac{\partial \mathcal{L}}{\partial y_0} = \mathbb{E}_0    & \left\{ U' (y_0) - \mu_0 ' \left. \frac{\partial f}{\partial y_t} \right|_{\left( y_1, y_0, y_{-1}, \epsilon_0 \right)} - \beta \mu_1 ' \left. \frac{\partial f}{\partial y_{t-1}} \right|_{ \left( y_2, y_1, y_{0}, \epsilon_1 \right) } \right\}                   \label{eq:ramsey_t0}
            \\
            0=\frac{\partial \mathcal{L}}{\partial y_t} = \mathbb{E}_0    & \left\{ U' (y_t) - \beta^{-1}\mu_{t-1} ' \left. \frac{\partial f}{\partial y_{t+1}} \right|_{\left( y_{t}, y_{t-1}, y_{t-2}, \epsilon_{t-1} \right)}\right.     \notag
            \\
                                                                          & \left. - \mu_t' \left. \frac{\partial f}{\partial y_{t}} \right|_{\left( y_{t+1}, y_{t}, y_{t-1}, \epsilon_{t} \right)} - \beta \mu_{t+1}' \left. \frac{\partial f}{\partial y_{t-1}} \right|_{\left( y_{t+2}, y_{t+1}, y_{t}, \epsilon_{t+1} \right)} \right\},~t>0 \label{eq:ramsey_t}
          \end{align}
    \item First period is special, rendering the problem non-time invariant (and time-inconsistent)!
    \item However: \eqref{eq:ramsey_t} reduces to \eqref{eq:ramsey_t0} with the initial Lagrange multiplier set to $\mu_{-1}=0$
  \end{itemize}

\end{frame}


\begin{frame}{Inspecting the structure}
  \begin{witemize}
    \item  With $\mu_{-1}=0$, we have a time-invariant non-linear system of rational expectation equations for all $t$:
    \begin{align}
      0= \frac{\partial \mathcal{L}}{\partial \mu_t} = \mathbb{E}_0 & f(y_{t+1}, y_{t}, y_{t-1}, \varepsilon_{t}),                                                                                                                                                                                                                         \\
      0=\frac{\partial \mathcal{L}}{\partial y_t} = \mathbb{E}_0    & \left\{ U' (y_t) - \beta^{-1}\mu_{t-1} ' \left. \frac{\partial f}{\partial y_{t+1}} \right|_{\left( y_{t}, y_{t-1}, y_{t-2}, \epsilon_{t-1} \right)}\right.     \notag
      \\
                                                                    & \left. - \mu_t' \left. \frac{\partial f}{\partial y_{t}} \right|_{\left( y_{t+1}, y_{t}, y_{t-1}, \epsilon_{t} \right)} - \beta \mu_{t+1}' \left. \frac{\partial f}{\partial y_{t-1}} \right|_{\left( y_{t+2}, y_{t+1}, y_{t}, \epsilon_{t+1} \right)} \right\} \: ,
    \end{align}
    given initial and transversality conditions
    \item System is amenable to both perfect foresight solver and k-order perturbation methods
  \end{witemize}

\end{frame}

\begin{frame}{Steady state}
  \begin{witemize}
    \item Perturbation of the non-linear system can be accurate because, in the absence of shocks, the system may converge to a well-defined steady state
    \item However: need to find a steady state for both $\mu$ and $y$ first
    \item The steady-state values $\overline{y}$ and $\overline{\mu}$ need to satisfy
    \begin{align}
       & f(\overline{y}, \overline{y}, \overline{y}, 0) = 0                                                                                                                                                                                                                                                                                                                                                                                     \\
       & U' \left(\overline{y} \right) - \overline{\mu} ' \left( \beta^{-1} \left. \frac{\partial f}{\partial y_{t+1}} \right|_{\left( \overline{y}, \overline{y}, \overline{y}, 0 \right)} + \left. \frac{\partial f}{\partial y_{t}} \right|_{\left( \overline{y}, \overline{y}, \overline{y}, 0 \right)} + \beta \left. \frac{\partial f}{\partial y_{t-1}} \right|_{\left( \overline{y}, \overline{y}, \overline{y}, 0 \right)} \right) = 0
    \end{align}
    \item That is a potentially large non-linear system to solve!
  \end{witemize}
\end{frame}

\begin{frame}{Using a conditional steady state}
  \begin{witemize}
    \item Problem can be simplified if a conditional steady state can be provided for $y$
    \item Take the planner instruments like the interest rate as given and compute the corresponding steady state of the private sector\\
    $\rightarrow$ implies a steady state candidate $\tilde y$
    \item Given $\tilde y$, we need to solve a linear equation to obtain $\tilde \mu$:
    \begin{align}
       & U' \left(\tilde{y} \right) - \tilde{\mu} ' \underbrace{\left( \beta^{-1} \left. \frac{\partial f}{\partial y_{t+1}} \right|_{\left( \tilde{y}, \tilde{y}, \tilde{y}, 0 \right)} + \left. \frac{\partial f}{\partial y_{t}} \right|_{\left( \tilde{y}, \tilde{y}, \tilde{y}, 0 \right)} + \beta \left. \frac{\partial f}{\partial y_{t-1}} \right|_{\left( \tilde{y}, \tilde{y}, \tilde{y}, 0 \right)} \right)}_M = U' \left(\tilde{y} \right) - \tilde{\mu} 'M= 0
    \end{align}
    \item Iterate on value of instruments until convergence of $\tilde y$ and $\tilde \mu$ to $\overline y$ and $\tilde \mu$\\
    $\rightarrow$ increases success probability of finding steady state to the problem
    \item However: steady state may not be unique\\
    $\rightarrow$ starting point for the instruments may matter
  \end{witemize}
\end{frame}


\begin{frame}{Welfare assessment}
  \begin{witemize}
    \item We are often interested in explicitly evaluating the welfare function
    \begin{equation}
      \max_{\left\{y_t\right\}_{t=0}^\infty} \mathbb{E}_0 \sum_{t=0}^\infty\beta^{t}U(y_t)  \: ,
    \end{equation}
    \item Welfare at time $t$ can be recursively written (making dependence on state variables $x_{t-1}$ and shocks $\varepsilon_t$ explicit):
    \begin{align}
      W\left( x_{t-1}, \varepsilon_t \right) = U\left( x_{t-1}, \varepsilon_t \right) + \beta \mathbb{E}_t W\left( x_{t}, \varepsilon_{t+1} \right)
    \end{align}
    \item Careful: welfare objective cannot be included in the \texttt{model} block as it will cause singularity
    \item \texttt{evaluate\_planner\_objective} allows evaluating this function
    \item Important: proper welfare assessment typically requires at least a second-order approximation of the model, unless there is an undistorted steady state \parencite{Woodford2002}
  \end{witemize}
\end{frame}


\begin{frame}{Conditional vs. unconditional welfare}
  \begin{itemize}
    \item \texttt{evaluate\_planner\_objective} reports two different welfare measures commonly of interest to policy-makers:
  \end{itemize}
          \begin{enumerate}
            \item Conditional welfare at time $t=0$:
                  \begin{align}
                    W \left( x_{-1}, \varepsilon_0, \right)
                  \end{align}

                  \begin{itemize}
                     \item takes into account the dynamics starting from the initial state $x_{-1}$ and the observed current shock $\varepsilon_0$
                     \item allows accounting for transition behavior from old policy regime to new optimal policy regime
                     \item Dynare reports value for two particular initial multipliers $\mu_{-1}$ states:
                           \begin{itemize}
                             \item $\mu_{-1} = \overline{\mu}$ (steady state multiplier)
                             \item $\mu_{-1} = 0$ (no initial commitment by planner)
                           \end{itemize}
                  \end{itemize}
            \item Unconditional welfare
                  \begin{align}
                    \mathbb{E}\left[ W\left( x_{t-1}, \varepsilon_t, \right) \right] = \frac{\mathbb{E}\left[U\left( x_{t-1}, \varepsilon_t, \right)\right]}{1-\beta}
                  \end{align}
                  $\rightarrow$ corresponds to the long-run expected value of welfare
          \end{enumerate}
\end{frame}


\subsection{Practice}

\begin{frame}{In Dynare}
  \begin{witemize}
    \item $U$ is stated in the \texttt{planner\_objective} command
    \item \texttt{ramsey\_constraints} block contains constraints, e.g. ZLB
    \item \texttt{ramsey\_model} command generates the FOC with options
          \begin{itemize}
            \item \texttt{planner\_discount}: central planner's discount factor
            \item \texttt{instruments}: variables for the computation of the steady state unde optimal policy\\
                  $\rightarrow$ \texttt{steady\_state\_model} block or \texttt{\_steadystate.m} file
          \end{itemize}
    \item \texttt{initval}: initial instruments values
    \item \texttt{histval}: initial state variables values
    \item \texttt{shocks}: specifies $\varepsilon_0$
    \item call \texttt{stoch\_simul} and \texttt{evaluate\_planner\_objective}
  \end{witemize}
\end{frame}


\begin{frame}{Dynare output}
  \begin{witemize}
    \item \texttt{oo\_.planner\_objective\_value} subfields:
          \begin{witemize}
            \item \texttt{unconditional} contains $\mathbb{E} [W]$
            \item \texttt{conditional} contains $W \left( x_{-1}, \varepsilon_0 \right)$ in subfields
                  \begin{itemize}
                    \item \texttt{steady\_initial\_multiplier} when $\mu_{-1} = \overline{\mu}$
                    \item \texttt{zero\_initial\_multiplier} when $\mu_{-1} = 0$\\
                          $\rightarrow$ optimal policy surprise for private agents!
                  \end{itemize}
          \end{witemize}
  \end{witemize}
\end{frame}


\begin{frame}{A New-Keynesian model}
  \begin{itemize}
    \item Consider model based on \textcite{CMR2007} \\
    $\rightarrow$ shipped with Dynare (\texttt{Ramsey\_Example.mod})
    \item Representative household chooses consumption $C_t$ and hours worked $h_t$ to maximize discounted utility
          \begin{equation}
            \mathbb{E}_0 \sum_{t=0}^\infty\beta^{t} \left[ \log \left( C_t \right) -\frac{\chi}{2} h_{t}^2\right]
          \end{equation}
    \item The budget constraint is given by
          \begin{equation}
            \frac{B_{t}}{P_t} =
            (1+r_{t-1})\frac{B_{t-1}}{P_t}-C_{t}+\frac{W_t}{P_t}h_{t} +\Theta_{t} \: ,
          \end{equation}
          where
          \begin{itemize}
            \item $B_t$: nominal private bonds (in zero net supply)
            \item $r_t$: nominal net interest rate
            \item $W_t$: nominal wage
            \item $\Theta_t$: firm profits
            \item $P_t$: price level
          \end{itemize}
  \end{itemize}
\end{frame}

\begin{frame}{Household FOCs}
  \begin{itemize}
    \item The consumption FOC is given by the familiar Euler equation
          \begin{equation}
            \frac{1}{C_{t}} =\beta \mathbb{E}_{t}\left \{\frac{1}{C_{t+1}}\frac{1+r_{t}}{\Pi_{t+1}}\right \} \: ,                                                                        \end{equation}
          where $\Pi_t$ is the inflation rate
    \item The labor FOC can be written as
          \begin{equation}\label{eq:labor}
            \chi h_tC_t=\frac{W_t}{P_t}
          \end{equation}
  \end{itemize}
\end{frame}


\begin{frame}{Firm setup}
  \begin{itemize}
    \item There is a continuum of monopolistically competitive firms $j$ with pricing power
    \item Production is linear in labor
          \begin{equation}
            C_{i,t}=e^{Z_t}h_{i,t} \: ,
          \end{equation}
          with TFP $Z_t$ following
          \begin{equation}
            Z_t = \rho Z_{t-1} + u_t \: , u \overset{iid}{\sim}\mathcal{N}(0,\sigma_z^2)
          \end{equation}
    \item Firm $j$ faces a downward-sloping demand curve for its variety
          \begin{equation}
            C_{j,t} = \left(\frac{P_{j,t}}{P_t}\right)^{-\theta} \: ,
          \end{equation}
          where $\theta$ is the Dixit-Stiglitz elasticity of substitution
    \item Price adjustment is subject to a \textcite{Rotemberg1982} quadratic cost
          \begin{equation}
            \frac{\phi}{2}\left(\frac{P_{j,t}}{P_{j,t-1}}-1\right)^2 C_{j,t}
          \end{equation}
  \end{itemize}
\end{frame}

\begin{frame}{Firm problem}
  \begin{itemize}
    \item Firm $j$ maximizes its after-tax profits
          \begin{equation}
            (1+\tau)\frac{P_{j,t}}{P_t}C_{j,t}-MC_t C_{j,t}-\frac{\phi}{2}\left(\frac{P_{j,t}}{P_{j,t-1}}-1\right)^2 C_{j,t} \: ,
          \end{equation}
          where marginal cost is common to all firms:
          \begin{equation}
            MC_t = \frac{W_t}{P_t e^{Z_t}} \overset{\eqref{eq:labor}}{=} \frac{\chi h_tC_t}{e^{Z_t}}
          \end{equation}
    \item In a symmetric equilibrium, the pricing FOC is the well-known New Keynesian Phillips curve
          \begin{equation}
            \left[\tau-\frac{1}{\theta-1}\right](1-\theta) +\theta\left(\frac{\chi
              h_t C_t}{e^{Z_t}}-1\right) = \phi(\Pi_t-1)\Pi_t - \beta \mathbb{E}_t\phi(\Pi_{t+1}-1)\Pi_{t+-1}
          \end{equation}
  \end{itemize}
\end{frame}

\begin{frame}{Private sector equilibrium}
  \begin{itemize}
    \item In equilibrium, the private sector FOCs given monetary policy can be represented as
          \begin{align}
            \frac{1}{C_{t}}            & =\beta \mathbb{E}_{t}\left \{ \frac{1}{C_{t+1}}\frac{1+r_{t}}{\Pi
            _{t+1}}\right \}                                                                                 \\
            \left[\tau-\frac{1}{\theta-1}\right](1-\theta) +\theta\left(\frac{\chi
            h_t C_t}{e^{Z_t}}-1\right) & = \phi(\Pi_t-1)\Pi_t - \beta \mathbb{E}_t\phi(\Pi_{t+1}-1)\Pi_{t+1} \\
            e^{Z_{t}}h_{t}             & = C_{t}\left[1+\frac{\phi }{2}\left( \Pi _{t}-1\right) ^{2}\right]  \\
            Z_t                        & = \rho Z_{t-1} + u_t
          \end{align}
  \end{itemize}

\end{frame}


\begin{frame}{Ramsey problem}
  \begin{itemize}
    \item The corresponding Ramsey problem is given by
          \begin{equation}
            \begin{split}
              \max \mathbb{E}_{0}\sum_{t=0}^\infty & \beta ^{t}\Bigg\{ \ln C_{t}-\frac{\chi}{2}h_t^2 -\mu_{1,t}\left( \frac{1}{C_{t}}- \beta \frac{1}{C_{t+1}}\frac{1+r_{t}}{\Pi
              _{t+1}}\right)                                                                                                                                                     \\
                                                   & -\mu_{2,t}\Bigg(  \left[\tau-\frac{1}{\theta-1}\right](1-\theta) +\theta\left(\frac{\chi
              h_t C_t}{e^{Z_t}}-1\right)                                                                                                                                         \\
                                                   & \quad\quad\quad- \phi(\Pi_t-1)\Pi_t + \beta \phi(\Pi_{t+1}-1)\Pi_{t-1}\Bigg)                                                \\
                                                   & -\mu _{3,t}\left(e^{Z_{t}}h_{t} - C_{t}\left[1+\frac{\phi }{2}\left( \Pi _{t}-1\right) ^{2}\right]\right)                   \\
                                                   & -\mu_{4,t}\left(Z_t - \rho Z_{t-1} - u_t\right)\Bigg\}
            \end{split}
          \end{equation}
  \end{itemize}
\end{frame}

\begin{frame}[fragile]{Dynare code}{\texttt{Ramsey\_Example\_commitment.mod}}
  \begin{itemize}
    \item We declare all 5 variables, including the interest rate set by the Ramsey planner:
  \end{itemize}
  \small
  \begin{minted}{c++}
var(log) C  $C$                 (long_name='Consumption')
    pi      $\pi$               (long_name='Gross inflation')
    h       $h$                 (long_name='hours worked');

var Z       $Z$                 (long_name='TFP')
    r       $r$                 (long_name='Net nominal interest rate');

varexo epsilon ${\varepsilon}$     (long_name='TFP shock');
parameters beta     ${\beta}$       (long_name='discount factor')
        theta       ${\theta}$      (long_name='substitution elasticity')
        tau         ${\tau}$        (long_name='labor subsidy')
        chi         ${\chi}$        (long_name='labor disutility')
        phi         ${\phi}$        (long_name='price adjustment costs')
        rho         ${\rho}$        (long_name='TFP autocorrelation');
\end{minted}
\end{frame}



% \begin{frame}[fragile]{Dynare code: \texttt{Ramsey\_Example\_commitment.mod}}
%   \begin{itemize}
%     \item We declare all 5 variables, including the interest rate set by the Ramsey planner:
%   \end{itemize}
%   \begin{minted}{c++}
% beta=0.99;
% theta=5;
% phi=100;
% rho=0.9;
% tau=1/(theta-1);
% chi=1;

% \end{minted}
% \end{frame}

\begin{frame}[fragile]{Dynare code (continued)}
  \begin{itemize}
    \item The model only features 4 equations for 5 variables:
  \end{itemize}
  \begin{minted}{c++}
model;
    [name='Euler equation']
    1/(1+r)=beta*C/(C(+1)*pi(+1));
    [name='Firm FOC']
    (tau-1/(theta-1))*(1-theta)+theta*(chi*h*C/(exp(Z))-1)=
        phi*(pi-1)*pi-beta*phi*(pi(+1)-1)*pi(+1);
    [name='Resource constraint']
    C*(1+phi/2*(pi-1)^2)=exp(Z)*h;
    [name='TFP process']
    Z=rho*Z(-1)+epsilon;
 end;
\end{minted}
\end{frame}

\begin{frame}[fragile]{Dynare code (continued)}
  \begin{itemize}
    \item We the steady state conditional on the interest rate $r$:
  \end{itemize}

  \begin{minted}{c++}
steady_state_model;
    Z=0;
    pi=(r+1)*beta;
    C=sqrt((1+1/theta*((1-beta)*(pi-1)*pi-(tau-1/(theta-1))*(1-theta)))/
        (chi*(1+phi/2*(pi-1)^2)));
    h=C*(1+phi/2*(pi-1)^2);
    y_nat=sqrt((theta-1)/theta*(1+tau)/chi);
    y_gap=log(C)-log(y_nat);
end;

shocks;
    var epsilon = 0.01^2;
end;
\end{minted}
\end{frame}

\begin{frame}[fragile]{Dynare code (continued)}
  Finally, we set the initial value and the planner objective:
  \begin{minted}{c++}
initval;
    r=1/beta-1;
end;
planner_objective log(C)-chi/2*h^2;

ramsey_model(instruments=(r),planner_discount=beta); 

//conduct stochastic simulations of the Ramsey problem
stoch_simul(order=2,irf=20,periods=0,irf_plot_threshold=0)
     LOG_pi LOG_h r LOG_C Z;
evaluate_planner_objective;     
\end{minted}
\end{frame}


\section{Discretionary policy}

\subsection{Theory}

\begin{frame}{Discretionary policy equilibrium}
  \begin{itemize}
    \item Planner re-optimizes each period and is expected to do so\\
          $\rightarrow$ no commitment!
    \item Denote with $y_t$ the endogenous variables and with $u_t$ the policy instruments
    \item Solution is only well-known for quadratic loss function
          \begin{align}
            L_0 = \sum_{t=0}^\infty \beta^{t}(y_t'Wy_t + u_t'Qu_t)
          \end{align}
          with linear constraints:
          \begin{equation}
            A_0 y_t = A_1y_{t-1} + A_2\mathbb{E}_t y_{t+1} + A_3
            u_t + A_4 \mathbb{E}_t u_{t+1} + A_5 \varepsilon_t \: ,
          \end{equation}
          where $\varepsilon_t$ are the structural shocks and $W,Q,A_i$ are coefficient matrices
  \end{itemize}
\end{frame}

\begin{frame}{Solution structure}
  \begin{itemize}
    \item We are looking for a Markov-perfect Stackelberg-Nash equilibrium:
          \begin{itemize}
            \item Stackelberg leader: today's planner
            \item Stackelberg followers: private sector agents and tomorrow's planner
          \end{itemize}
    \item We conjecture our solution will be linear:
          \begin{align}
            y_t & = H_1 y_{t-1} + H_2 \varepsilon_t \\
            u_t & = F_1 y_{t-1} + F_2 \varepsilon_t \: ,
          \end{align}
          where $H_1,\;H_2,\;F_1,\;F_2$ are unknown matrices to be determined
    \item Plugging the guess into the constraints allows deriving a matrix equation whose fixed point corresponds to the solution
  \end{itemize}
\end{frame}

\begin{frame}{\textcite{Dennis2007} Solution algorithm}
  {\scriptsize
    \begin{description}
      \item[Step 1] For $i=0$, initialize $H_1^{(i)}$, $H_2^{(i)}$, $F_1^{(i)}$,
            $F_2^{(i)}$
      \item[Step 2] Compute $D^{(i)} = A_0 - A_2 H_1^{(i)} - A_4 F_1^{(i)}$
      \item[Step 3] Solve the so-called Sylvester equation
            \begin{equation}
              P^{(i)} = W + \beta {F_1^{(i)}}'Q F_1^{(i)} + \beta H_1^{(i)} P^{(i)} H_1^{(i)}
            \end{equation}
      \item[Step 4] Set
            \begin{align*}
              F_1^{(i+1)} & =  -(Q+A_3'(                                            %{D^{(i)}}')^{-1}P^{(i)}
              {D^{(i)}}^{-1}A_3)^{-1}A_3'({D^{(i)}}')^{-1}P^{(i)}{D^{(i)}}^{-1}A_1, \\
              F_2^{(i+1)} & =   -(Q+A_3'({D^{(i)}}')^{-1}P^{(i)}
              {D^{(i)}}^{-1}A_3)^{-1}A_3'({D^{(i)}}')^{-1}P^{(i)}{D^{(i)}}^{-1}A_5, \\
              H^{(i+1)}_1 & = {D^{(i)}}^{-1}(A_1+A_3F_1^{(i+1)}),                   \\
              H^{(i+1)}_2 & = {D^{(i)}}^{-1}(A_5+A_3F_2^{(i+1)}).                   \\
            \end{align*}
      \item[Step 5] Go back to Step 2, until convergence.
    \end{description}
  }
\end{frame}

\subsection{Practice}

\begin{frame}{In Dynare}
  \begin{witemize}
    \item Quadratic loss function must be provided to Dynare in \texttt{planner\_objective}
    \item Careful: no linear terms are allowed in the loss function and the variables in the loss function must be mean 0 
    \item Instrument name and full steady state must be provided (not a conditional one)
    \item \texttt{discretionary\_policy} to compute discretionary policy, generate IRFs, etc\\
    $\rightarrow$ similar to \texttt{ramsey\_model} followed by \texttt{stoch\_simul}
  \end{witemize}
\end{frame}

\begin{frame}[fragile]{Dynare code}{\texttt{Ramsey\_Example\_discretionary.mod}}
  \begin{itemize}
    \item We need to use a quadratic loss function derived from the nonlinear planner objective:
  \end{itemize}
\small
\begin{minted}{c++}
 steady_state_model;
    r=1/beta-1;
    Z=0;
    pi=(r+1)*beta;
    C=sqrt((1+1/theta*((1-beta)*(pi-1)*pi-(tau-1/(theta-1))*(1-theta)))/
    (chi*(1+phi/2*(pi-1)^2)));
    h=C*(1+phi/2*(pi-1)^2);
    y_nat=sqrt((theta-1)/theta*(1+tau)/chi);
    y_gap=log(C)-log(y_nat);
end;

planner_objective pi_hat^2*theta/((theta-1)/phi) + y_gap^2;

discretionary_policy(instruments=(r),planner_discount=beta); 
\end{minted}
\end{frame}

\section{Optimal simple rules}

\subsection{Theory}

\begin{frame}{Program}
  \begin{itemize}
    \item The general OSR problem involves the planner picking the optimal parameters $\gamma$ of a policy rule
          \begin{equation}
            u_t = g(y_{t+1},y_{t},y_{t-1},\varepsilon_t, \gamma) \: ,
          \end{equation}
          where $u_t$ are the policy instruments
          to maximize the objective function
          \begin{equation}
            \max_{\gamma} \mathbb{E}_0 \sum_{t=0}^\infty\beta^{t}U(y_t) 
          \end{equation}
    \item Optimization is conducted subject to the private sector equilibrium conditions:
          \begin{equation}
            \mathbb{E}_{t}f(y_{t+1}, y_{t}, y_{t-1}, \varepsilon_{t})=0,~\forall \:  t \geq 0 \: ,
          \end{equation}
          taking $y_{-1}$ as given and imposing suitable transversality conditions
  \end{itemize}

\end{frame}

\begin{frame}{Program}
  \begin{witemize}
    \item Up to Dynare 6, only special linear-quadratic (LQ) versions of the problem could be solved:
          \begin{align}
                         & \min_\gamma \mathbb{E} \left[ y_t'W y_t \right]                        \\
            \text{s.t.}~ & A_1 \mathbb{E}_t y_{t+1} + A_2 y_t + A_3 y_{t-1} + A_4 \varepsilon_t = 0 \: ,
          \end{align}
          where $A_i$ are coefficient matrices obtained from using first-order perturbation
    \item $W_t$ is the weighting matrix penalizing variances and covariances of the variables\\
          $\rightarrow$ specified in Dynare using the \texttt{optim\_weights} block\\
    \item Loss function cannot be automatically derived from \texttt{planner\_objective}
    \item Disadvantage: LQ setup restricted to first-order approximation, which is often insufficient for proper welfare evaluation (and comparisons to Ramsey)
    \item Advantage: LQ setup makes e.g. using analytical derivatives possible\\ (\texttt{analytical\_derivation}-option)
  \end{witemize}
\end{frame}


\begin{frame}{Legacy implementation in Dynare}
  \begin{witemize}
    \item \texttt{osr\_params} block specifies parameters $\gamma$ to optimize over
    \item \texttt{optim\_weights} block specifies weighting matrix $W$
    \item \texttt{osr\_param\_bounds} block sets bounds for parameter optimization
    \item  \texttt{osr} command computes solution and results (analogously to \texttt{stoch\_simul})
    \item  output and metadata located in fields \texttt{M\_.osr} and \texttt{oo\_.osr}
  \end{witemize}
\end{frame}


\begin{frame}{New implementation in Dynare 7}
  \begin{witemize}
    \item Starting with Dynare 7 (and already available in the unstable version), a new implementation of the OSR algorithm is available
    \item Instead of \texttt{optim\_weights} block a \texttt{planner\_objective} command can be specified
    \item  \texttt{osr} command computes solution at \texttt{order}$\geq 1$
    \item \texttt{evaluate\_planner\_objective} command computes welfare at the solution
  \end{witemize}
\end{frame}



\subsection{Practice}

\begin{frame}[fragile]{Dynare code}{\texttt{Ramsey\_Example\_OSR.mod}}
  \begin{itemize}
    \item Need to include the Taylor rule whose coefficient we optimize:
    \begin{equation}
      r-\bar r =\alpha (\pi-\pi^*)
    \end{equation}
  \end{itemize}\vspace{-0.5cm}
  \small
  \begin{minted}{c++}
  model;
      [name='Euler equation']
      1/(1+r)=beta*C/(C(+1)*pi(+1));
      [name='Firm FOC']
      (tau-1/(theta-1))*(1-theta)+theta*(chi*h*C/(exp(Z))-1)
      =phi*(pi-1)*pi-beta*phi*(pi(+1)-1)*pi(+1);
      [name='Resource constraint']
      C*(1+phi/2*(pi-1)^2)=exp(Z)*h;
      [name='TFP process']
      Z=rho*Z(-1)+epsilon;
      [name='Taylor rule']
      r=pi_star/beta-1+alpha*(pi-pi_star);
  end;
\end{minted}
\end{frame}



\begin{frame}[fragile]{Dynare code}{\texttt{Ramsey\_Example\_OSR.mod}}
  \begin{minted}{c++}
planner_objective log(C)-chi/2*h^2;

//define OSR parameters to be optimized
osr_params alpha;
//starting value for OSR parameter
alpha = 1.5;
//define bounds for OSR during optimization
osr_params_bounds;
    alpha, 0, 100;
end;
//compute OSR and provide output
osr(opt_algo=9,order=2,planner_discount=beta,irf_plot_threshold=0) 
  LOG_pi LOG_h r LOG_C Z;
evaluate_planner_objective;
\end{minted}
\end{frame}

% \section{Conclusion}

% \begin{frame}{Conclusion}
%   \begin{itemize}
%     \item Manual
%     \item \texttt{tests} subfolder: see subfolders
%           \begin{itemize}
%             \item \texttt{optimal\_policy}
%             \item \texttt{discretionary\_policy}
%           \end{itemize}
%           for examples
%     \item Pruning for Ramsey optimal policy is on the stack
%   \end{itemize}
% \end{frame}


\begin{frame}[plain,allowframebreaks]{References}
  \printbibliography
\end{frame}

\end{document}
